\section{Mathematical visualization atlas}
\label{sec:visualatlas}

This atlas supports algorithm explanations with actual data where available,
and with explicitly labelled reduced examples elsewhere. Plots are generated
by \texttt{make plots-math} and embedded only when the corresponding PDF
exists. Every figure has an SVG, PNG and PDF counterpart under
\texttt{build/math\_plots/}. The complete topic-to-file map is recorded in
\texttt{docs/visualization\_index.md}. None of these figures is a
measurement of algorithm runtime.

\newcommand{\EulerChart}[3]{%
 \begin{figure}[p]\centering
 \IfFileExists{../../build/math_plots/#1.pdf}{%
  \includegraphics[width=0.92\linewidth,height=0.37\textheight,keepaspectratio]{../../build/math_plots/#1.pdf}%
 }{\fbox{\parbox{0.76\linewidth}{%
  Figure not generated: \texttt{\detokenize{#1.pdf}}. Run \texttt{make plots-math}
  before rebuilding this report.}}}%
 \caption{Project Euler #2: #3}\label{fig:#1}
 \end{figure}
}

\newcommand{\MathChart}[2]{%
 \begin{figure}[p]\centering
 \IfFileExists{../../build/math_plots/#1.pdf}{%
  \includegraphics[width=0.92\linewidth,height=0.37\textheight,keepaspectratio]{../../build/math_plots/#1.pdf}%
 }{\fbox{\parbox{0.76\linewidth}{%
  Figure not generated: \texttt{\detokenize{#1.pdf}}. Run \texttt{make plots-math}
  before rebuilding this report.}}}%
 \caption{#2}\label{fig:#1}
 \end{figure}
}

\subsection{Integer geometry and spirals}
\EulerChart{pythagorean_perimeters}{9, 39, 75}{Euclid's formula enumerates
right-triangle perimeter counts; the dashed marker gives the full problem
39 winner $p=840$. Singular-count growth is sampled only to $p=1500$.}
\EulerChart{spiral_diagonal_sum}{28}{Sum of the four corners of each odd-width
ring through the full $1001\times1001$ target.}
\EulerChart{spiral_prime_ratio}{58}{Prime density on spiral diagonals sampled
through side length 2501, with a 10-percent reference line.}
\EulerChart{coordinate_right_triangles}{91}{Lattice-coordinate geometry
on a small grid; the full counting problem is not implemented here.}
\EulerChart{almost_equilateral}{94}{The $5,5,6$ example has an integral area
and differs from an equilateral triangle by one unit.}

\subsection{Grids and dynamic programming}
\EulerChart{grid_largest_product}{11}{Actual $20\times20$ input grid
from the repository. The highlighted window produces the verified
maximum product $70,600,674$.}
\EulerChart{lattice_paths}{15}{Binomial coefficients and a sample
six-by-six lattice path. The chart is not the full problem input.}
\EulerChart{maximum_path_sum_demo}{18, 67}{A reduced four-row triangle
shows how bottom-up dynamic programming identifies the maximum path sum.}
\EulerChart{path_sum_directions}{81, 82, 83}{Three shortest-path
models share one five-by-five example matrix but allow different moves.
The full-sized input matrices are not supplied.}
\EulerChart{sudoku_demo}{96}{A solved nine-by-nine example, distinguishing
givens from cells filled by backtracking.}
\EulerChart{triangle_containment}{102}{The signs of three oriented cross
products decide whether the origin is inside a triangle.}

\subsection{Discrete counting and networks}
\EulerChart{magic_5gon}{68}{One valid five-gon configuration whose
five triples have the same sum.}
\EulerChart{monopoly_dice_kernel}{84}{The probability mass function of
two four-sided dice. This is only the movement kernel, not Monopoly's
complete Markov model with jail and card rules.}
\EulerChart{counting_rectangles}{85}{The number of rectangular
subgrids across width and height, with a target contour at two million.}
\EulerChart{cuboid_routes}{86}{Integer-length unfolded surface
routes for a restricted set of cuboids.}
\EulerChart{cube_digit_pairs}{90}{A valid digit placement on two cubes,
where six and nine can be exchanged by rotating the cube.}
\EulerChart{minimal_network}{107}{Kruskal's minimum-spanning-tree
procedure, shown on an explicitly synthetic graph.}
\EulerChart{darts_checkout_counts}{109}{Distinct checkout
combinations under score 80, with an obligatory final double.}
\EulerChart{block_tilings}{114, 115, 116, 117}{Four exact
dynamic-programming recurrences on short rows; the second curve is a
scaled illustration, not the full problem-115 threshold.}

\subsection{Pascal's pyramid}
\MathChart{pascal_pyramid_tetrahedron}{Every trinomial coefficient through
layer 14. The integer coordinates $(i,j,k)$ identify the coefficient of
$x^iy^jz^k$ and fill a tetrahedron bounded by $i+j+k\leq14$.}
\MathChart{pascal_pyramid_layer_surface}{Layer 24 drawn in barycentric
coordinates. Height and color encode the base-10 logarithm of the exact
coefficient, so the three permutation-symmetric edges remain visible.}
\MathChart{pascal_pyramid_parity}{Odd coefficients through layer 31. Removing
the even coefficients exposes the self-similar mod-2 tetrahedron.}
\clearpage
